\documentclass[11pt]{amsart}

\usepackage[T1]{fontenc}
\usepackage{lmodern}
\usepackage{amsmath,amssymb,amsthm,mathtools}
\usepackage[margin=1.1in]{geometry}
\usepackage[hidelinks]{hyperref}

\newtheorem{theorem}{Theorem}
\newtheorem{proposition}[theorem]{Proposition}
\newtheorem{corollary}[theorem]{Corollary}
\newtheorem{remark}[theorem]{Remark}

\newcommand{\T}{\mathbb T}
\newcommand{\Z}{\mathbb Z}
\newcommand{\R}{\mathbb R}
\newcommand{\C}{\mathbb C}
\newcommand{\conv}{\operatorname{conv}}
\newcommand{\supp}{\operatorname{supp}}
\newcommand{\Span}{\operatorname{span}}

\title[Non-invariance of finite Fourier coordinate subspaces]
{Non-invariance of finite Fourier coordinate subspaces\\
under cubic--quintic Schr\"odinger evolution}
\date{October 26, 2026}

\begin{document}
\begin{abstract}
For the unprojected fourth-order cubic--quintic Schr\"odinger flow on a flat
periodic torus, we prove that no finite Fourier coordinate subspace with at
least two frequencies is invariant.  In fact, every trigonometric polynomial
with at least two nonzero Fourier coefficients has an initially absent Fourier
mode with nonzero time derivative.  The proof isolates an extremal quintic
interaction that cannot be canceled by the cubic nonlinearity or the diagonal
linear propagator.  We also classify the invariant finite Fourier coordinate
subspaces and distinguish the result from claims about approximate Galerkin
models or general nonlinear invariant manifolds.
\end{abstract}
\maketitle

\section{Equation and statement}

Let $d\ge 1$ and $\T^d=(\R/2\pi\Z)^d$.  Consider
\begin{equation}\label{eq:flow}
 i\partial_t u=(\Delta+\kappa^2)^2u-\beta|u|^2u+\gamma|u|^4u,
 \qquad (t,x)\in\R\times\T^d,
\end{equation}
where $\kappa>0$, $\beta\in\R$, and $\gamma\in\R\setminus\{0\}$.
In particular, the claims apply to the focusing-cubic/defocusing-quintic
case $\beta,\gamma>0$.

For a finite set $S\subset\Z^d$, define its Fourier coordinate subspace
\begin{equation}\label{eq:ES}
 E_S:=\Span_{\C}\{e^{in\cdot x}:n\in S\}.
\end{equation}
Invariance means that the solution of \eqref{eq:flow} remains in $E_S$
whenever its initial datum lies in $E_S$, for as long as that solution exists.

\begin{theorem}[No finite multimode Fourier coordinate subspace is invariant]
\label{thm:noninvariance}
If $S\subset\Z^d$ is finite and $|S|\ge 2$, then $E_S$ is not
invariant under \eqref{eq:flow}.  More precisely, there exist
$u_0\in E_S$ and $r\in\Z^d\setminus S$ such that the corresponding
smooth local solution satisfies
\begin{equation}\label{eq:derivative-nonzero}
 \widehat u(r,0)=0,
 \qquad \partial_t\widehat u(r,0)\ne 0.
\end{equation}
Consequently, $u(t)\notin E_S$ for every sufficiently small $t\ne0$.
\end{theorem}

\begin{proof}
Choose an extreme point $p$ of the compact polytope $\conv(S)$.
Every extreme point of the convex hull of a finite set belongs to that set,
so $p\in S$.  Choose $q\in S$ with $q\ne p$ and set
\begin{equation}\label{eq:r}
 r:=3p-2q\in\Z^d.
\end{equation}
If $r\in S$, then
\[
 p=\frac23q+\frac13r
\]
would express the extreme point $p$ as a nontrivial convex combination
of two distinct elements of $\conv(S)$.  Therefore $r\notin S$.

Take $a,b\in\C\setminus\{0\}$ and put
\begin{equation}\label{eq:two-mode}
 u_0(x):=a e^{ip\cdot x}+b e^{iq\cdot x}\in E_S.
\end{equation}
The linear operator $L:=(\Delta+\kappa^2)^2$ is diagonal in the
Fourier basis:
\[
 L e^{in\cdot x}=(\kappa^2-|n|^2)^2e^{in\cdot x}.
\]
Thus $\widehat{Lu_0}(r)=0$.

Next,
\[
 |u_0|^2u_0=u_0^2\overline{u_0}
\]
has Fourier support contained in
$\{p,q,2p-q,2q-p\}$.  Since $p\ne q$, none of these
frequencies equals $r=3p-2q$; hence
\begin{equation}\label{eq:cubic-zero}
 \widehat{|u_0|^2u_0}(r)=0.
\end{equation}

Finally,
\[
 |u_0|^4u_0=u_0^3\overline{u_0}^{\,2}.
\]
A term in this product has frequency
\[
 jp+(3-j)q-\ell p-(2-\ell)q
   =(j-\ell)p+(1-j+\ell)q,
\]
where $j\in\{0,1,2,3\}$ and $\ell\in\{0,1,2\}$.
To obtain $r=3p-2q$, we must have
\[
 (j-\ell-3)(p-q)=0.
\]
Because $p\ne q$, this implies $j-\ell=3$, whose only allowed
solution is $(j,\ell)=(3,0)$.  Thus there is exactly one quintic
contribution at frequency $r$, namely
\begin{equation}\label{eq:quintic-coeff}
 \widehat{|u_0|^4u_0}(r)=a^3\overline b^{\,2}\ne0.
\end{equation}
Taking the $r$-th Fourier coefficient of \eqref{eq:flow} at $t=0$
and using \eqref{eq:cubic-zero}--\eqref{eq:quintic-coeff}, we get
\begin{equation}\label{eq:exact-derivative}
 \partial_t\widehat u(r,0)=-i\gamma a^3\overline b^{\,2}\ne0.
\end{equation}
Smooth trigonometric-polynomial data generate a smooth local solution
of the polynomial semilinear Schr\"odinger equation; in particular,
each Fourier coefficient is differentiable at $t=0$.  Since
$\widehat u(r,0)=0$, equation \eqref{eq:exact-derivative} gives
\[
 \widehat u(r,t)=-i\gamma a^3\overline b^{\,2}t+o(t)
 \qquad (t\to0).
\]
Its coefficient is nonzero for every sufficiently small $t\ne0$.
As $r\notin S$, the solution has left $E_S$.
\end{proof}

\section{Leakage from every finite multimode initial support}

The preceding theorem constructs a particular initial datum in each
coordinate subspace.  The following stronger statement applies to
\emph{every} finite Fourier superposition with at least two active modes.

\begin{theorem}[Universal instantaneous leakage for trigonometric polynomials]
\label{thm:all-data}
Let
\begin{equation}\label{eq:general-u}
 u_0(x)=\sum_{n\in T}c_n e^{in\cdot x},
 \qquad T\subset\Z^d\text{ finite},\quad |T|\ge2,
 \quad c_n\ne0\ (n\in T).
\end{equation}
There is a frequency $r\notin T$ such that
\[
 \partial_t\widehat u(r,0)\ne0.
\]
Consequently, the solution leaves $E_T$ for all sufficiently small
nonzero times.
\end{theorem}

\begin{proof}
Choose a real linear functional $\ell:\R^d\to\R$ taking distinct
values on $T$.  This is possible because one only needs to avoid the
finitely many hyperplanes
$\{v\in\R^d:v\cdot(n-m)=0\}$ for distinct $n,m\in T$.
Let $p$ and $q$ be the unique maximum and minimum points of $\ell|_T$,
respectively.  Then
\[
 \ell(p)>\ell(q),\qquad
 \ell(q)\le\ell(n)\le\ell(p)\quad(n\in T),
\]
with the appropriate equality occurring only at the relevant endpoint.
Set $r=3p-2q$.  Since
\[
 \ell(r)=3\ell(p)-2\ell(q)>\ell(p),
\]
we have $r\notin T$.

The exact quintic Fourier convolution is
\begin{equation}\label{eq:conv}
 \widehat{|u_0|^4u_0}(r)
 =\sum_{\substack{n_1,n_2,n_3,m_1,m_2\in T\\
                n_1+n_2+n_3-m_1-m_2=r}}
 c_{n_1}c_{n_2}c_{n_3}
 \overline{c_{m_1}}\,\overline{c_{m_2}}.
\end{equation}
Every tuple in this sum obeys
\[
 \begin{aligned}
 \ell(r)
 &=\ell(n_1)+\ell(n_2)+\ell(n_3)-\ell(m_1)-\ell(m_2)\\
 &\le 3\ell(p)-2\ell(q)=\ell(r).
 \end{aligned}
\]
Equality forces $n_1=n_2=n_3=p$ and $m_1=m_2=q$, because the
extrema of $\ell$ on $T$ are unique.  Hence the sum
\eqref{eq:conv} contains exactly one term:
\begin{equation}\label{eq:universal-quintic}
 \widehat{|u_0|^4u_0}(r)=c_p^3\overline{c_q}^{\,2}\ne0.
\end{equation}

By contrast, every frequency $s$ in the cubic convolution
$u_0^2\overline{u_0}$ satisfies
\[
 \ell(s)\le2\ell(p)-\ell(q)
  <3\ell(p)-2\ell(q)=\ell(r).
\]
It follows that $\widehat{|u_0|^2u_0}(r)=0$.
Since $L$ is diagonal and $r\notin T$, also
$\widehat{Lu_0}(r)=0$.  Equation \eqref{eq:flow} now yields
\begin{equation}\label{eq:universal-leakage}
 \partial_t\widehat u(r,0)
  =-i\gamma c_p^3\overline{c_q}^{\,2}\ne0.
\end{equation}
The local-in-time conclusion follows as in Theorem~\ref{thm:noninvariance}.
\end{proof}

\section{Classification and explicit example}

\begin{corollary}[Complete classification of finite Fourier coordinate subspaces]
\label{cor:classification}
For every finite $S\subset\Z^d$,
\[
 \boxed{\ E_S\text{ is invariant under \eqref{eq:flow}}
 \quad\Longleftrightarrow\quad |S|\le1.\ }
\]
\end{corollary}

\begin{proof}
Theorem~\ref{thm:noninvariance} excludes $|S|\ge2$.
The empty subspace $E_\varnothing=\{0\}$ is invariant.
For $S=\{p\}$, substitute $u(t,x)=A(t)e^{ip\cdot x}$ into
\eqref{eq:flow}.  The amplitude solves
\[
 iA'(t)=\big((\kappa^2-|p|^2)^2-\beta|A(t)|^2
                    +\gamma|A(t)|^4\big)A(t).
\]
Its modulus is constant, so
\[
 A(t)=A(0)e^{-i\Omega t},\qquad
 \Omega=(\kappa^2-|p|^2)^2-\beta|A(0)|^2+\gamma|A(0)|^4.
\]
Thus $E_{\{p\}}$ is invariant.
\end{proof}

\begin{proposition}[Fifth-harmonic generation from a cosine]
\label{prop:cos}
In one spatial dimension, let $u_0(x)=\cos x$.  For the solution
of \eqref{eq:flow},
\[
 \widehat u(5,0)=0,
 \qquad
 \widehat u(5,t)=-\frac{i\gamma}{32}t+o(t)
 \quad(t\to0).
\]
\end{proposition}

\begin{proof}
Apply Theorem~\ref{thm:noninvariance} with
$p=1$, $q=-1$, and $a=b=1/2$; then $r=5$ and
$a^3\overline b^{\,2}=1/32$.
\end{proof}

\section{Scope and interpretation}

\begin{remark}[Unprojected versus Galerkin-projected evolution]
For a finite set $S$, the Galerkin equation
\[
 i\partial_tu=P_S\big[(\Delta+\kappa^2)^2u
  -\beta|u|^2u+\gamma|u|^4u\big]
\]
remains in $E_S$ by construction.  This does not establish invariance
of $E_S$ for the original unprojected equation: $P_S$ discards the
newly generated harmonics and changes the dynamics.
\end{remark}

\begin{remark}[What the result does not prove]
Theorems~\ref{thm:noninvariance} and \ref{thm:all-data}
concern finite \emph{Fourier coordinate subspaces}.  They do not exclude
other finite-dimensional nonlinear invariant manifolds, rule out useful
finite-time approximation by Galerkin or envelope models, or establish
that every solution populates infinitely many frequencies.  They also
do not derive the cubic--quintic equation from a proposed higher-derivative
parent theory.  A controlled reduction must separately justify excluded
modes and quantify its approximation error.
\end{remark}

\end{document}
